\setchoices{label=(\Alph*)}

\begin{question}[index=1]
曲线 $y=(x+2)\mathrm e^{\frac{1}{x}}$ 的渐近线的条数为
\begin{choices}[columns=4]
\item $0$．
\item $1$．
\item $2$．
\item $3$．
\end{choices}
\end{question}
\QuestionBreak

\begin{question}[index=2]
设函数 $f(x)$ 可导，$F(x)=f(x)[1-|\ln(1+x)|]$，则
\begin{choices}[columns=1]
\item $F(x)$ 在 $x=0$ 处不可导．
\item 当 $F(x)$ 在 $x=0$ 处可导时，$f(0)=0$．
\item 当 $F(x)$ 在 $x=0$ 处可导时，$f(0)=1$．
\item $F'(0)=f'(0)$．
\end{choices}
\end{question}
\QuestionBreak

\begin{question}[index=3]
设 $L$ 为曲线 $\left(x-\dfrac{1}{2}\right)^2+\left(y-\dfrac{1}{2}\right)^2=\dfrac{1}{2}$，
取逆时针方向，
\[
\begin{aligned}
I&=\oint_L 4y\dd x+(x+y)^2\dd y,\\
J&=\oint_L 4x\dd x+(x+y)^2\dd y,\qquad
K=\oint_L 4xy\dd x+(x+y)^2\dd y,
\end{aligned}
\]
则 $I,J,K$ 的大小关系为
\begin{choices}[columns=2]
\item $I<K<J$．
\item $J<K<I$．
\item $I<J<K$．
\item $K<I<J$．
\end{choices}
\end{question}
\QuestionBreak

\begin{question}[index=4]
设 $u=f(x^2,y^2,xyz)$，函数 $z=g(x,y)$ 由方程
$\displaystyle\int_{\mathrm e^{xy}}^z\varphi(\mathrm e^{xy}+z-t)\dd t=xyz$
确定，其中 $f$ 可微，$\varphi$ 连续，且满足 $\varphi(z)-xy\ne0$，
则 $x\dfrac{\partial u}{\partial x}-y\dfrac{\partial u}{\partial y}=$
\begin{choices}[columns=2]
\item $2xf'_1+2yf'_2$．
\item $2xf'_1-2yf'_2$．
\item $2x^2f'_1+2y^2f'_2$．
\item $2x^2f'_1-2y^2f'_2$．
\end{choices}
\end{question}
\QuestionBreak

\begin{question}[index=5]
设 $\symbfit{A}=\begin{pmatrix}1&0&1\\5&4&9\\2&4&6\end{pmatrix}$，
$\symbfit{A}\symbfit{x}=\symbf{0}$ 的全部解构成的向量空间称为 $\symbfit{A}$ 的零空间，\\
$\symbfit{A}\symbfit{x}=\symbf{0}$ 的\mbox{基础解系}所含向量的个数称为 $\symbfit{A}$ 的零空间维数，则
\begin{choices}[columns=1]
\item $\symbfit{A}$ 的零空间为 $k(1,1,-1)^{\mathrm T}$，$k$ 为任意常数．
\item $\symbfit{A}$ 的零空间维数为 $2$．
\item $\symbfit{A}^{\mathrm T}$ 的零空间为 $k(1,1,-1)^{\mathrm T}+l(-1,1,0)^{\mathrm T}$，$k,l$ 为任意常数．
\item $\symbfit{A}^{\mathrm T}$ 的零空间维数为 $2$．
\end{choices}
\end{question}
\QuestionBreak

\begin{question}[index=6]
设 $\symbfit{A}$ 为 $3$ 阶非零矩阵，$\symbfit{A}^*$ 是 $\symbfit{A}$ 的伴随矩阵，
若 $\symbfit{A}^*+2\symbfit{A}=\symbfit{O}$，则 $\symbfit{A}^{10}=$
\begin{choices}[columns=2,label-pos=left]
\item $\begin{pmatrix}-4^5&&\\&-4^5&\\&&-4^5\end{pmatrix}$．
\item $\begin{pmatrix}-4^5&&\\&-4^5&\\&&4^5\end{pmatrix}$．
\item $\begin{pmatrix}-4^5&&\\&4^5&\\&&4^5\end{pmatrix}$．
\item $\begin{pmatrix}4^5&&\\&4^5&\\&&4^5\end{pmatrix}$．
\end{choices}
\end{question}
\QuestionBreak

\begin{question}[index=7]
设 $3$ 阶矩阵 $\symbfit{A},\symbfit{B}$ 满足
$\symbfit{A}^2+\symbfit{B}^2=\symbfit{A}\symbfit{B}+\symbfit{B}\symbfit{A}$，
且 $\symbfit{A}\ne\symbfit{B}$，则下列结论错误的是
\begin{choices}[columns=2]
\item $\symbfit{A}-\symbfit{B}$ 的秩为 $1$．
\item $\symbfit{A}-\symbfit{B}$ 的迹为 $0$．
\item $\symbfit{A}-\symbfit{B}$ 可有非零特征值．
\item $\symbfit{A}-\symbfit{B}$ 不可相似对角化．
\end{choices}
\end{question}
\QuestionBreak

\begin{question}[index=8]
设 $X_1,X_2,\ldots,X_{10}$ 是来自总体 $N(0,\sigma^2)$（$\sigma>0$）的简单随机样本，记
\[
\overline X=\frac{1}{10}\sum_{i=1}^{10}X_i,\qquad
S_1^2=\frac{1}{9}\sum_{i=1}^{10}(X_i-\overline X)^2,\qquad
S_2^2=\frac{1}{10}\sum_{i=1}^{10}X_i^2,\quad\text{则}
\]
\begin{choices}[columns=2,label-pos=left]
\item $\dfrac{\sqrt{10}\,\overline X}{S_1}\dist t(10)$．
\item $\dfrac{\sqrt{10}\,\overline X}{S_2}\dist t(10)$．
\item $\dfrac{9S_1^2}{\sigma^2}\dist\chi^2(9)$．
\item $\dfrac{10S_2^2}{\sigma^2}\dist\chi^2(9)$．
\end{choices}
\end{question}
\QuestionBreak

\begin{question}[index=9]
设 $X$ 与 $Y$ 为具有二阶矩的随机变量，且 $DX\ne0$，
若记 $G(a,b)=E\bigl\{[Y-(a+bX)]^2\bigr\}$，则 $G(a,b)$ 在驻点处的值为
\begin{choices}[columns=2]
\item $(1-\rho_{XY}^2)DY$．
\item $(1+\rho_{XY}^2)DY$．
\item $(1-\rho_{XY}^2)\operatorname{Cov}(X,Y)$．
\item $(1+\rho_{XY}^2)\operatorname{Cov}(X,Y)$．
\end{choices}
\end{question}
\QuestionBreak

\begin{question}[index=10]
设总体 $X$ 的分布为
\[
\begin{gathered}
P\{X=k\}=\theta^{k-1}(1-\theta),\quad k=1,2,\ldots,r,\\
P\{X=r+1\}=\theta^r,\quad \theta\in(0,1).
\end{gathered}
\]
从总体中抽取容量为 $n$ 的样本 $x_1,x_2,\ldots,x_n$，
其中有 $M$ 个取值为 $r+1$，\\
则 $\theta$ 的最大似然估计值为
\begin{choices}[columns=2,label-pos=left]
\item $\hat\theta=\dfrac{\displaystyle\sum_{i=1}^n x_i+n}{\displaystyle\sum_{i=1}^n x_i+M}$．
\item $\hat\theta=\dfrac{\displaystyle\sum_{i=1}^n x_i-n}{\displaystyle\sum_{i=1}^n x_i-M}$．
\item $\hat\theta=\dfrac{\displaystyle\sum_{i=1}^n x_i^2+n}{\displaystyle\sum_{i=1}^n x_i^2+M}$．
\item $\hat\theta=\dfrac{\displaystyle\sum_{i=1}^n x_i^2-n}{\displaystyle\sum_{i=1}^n x_i^2-M}$．
\end{choices}
\end{question}
\QuestionBreak
